\par\bigskip\SourceLeaf{163}
\Needspace{8\baselineskip}\pdfbookmark[1]{III. Commissionsbericht}{witwen-III}
\begin{center}
[III.]\par\medskip
[\textit{Commissionsbericht.}]
\end{center}
\noindent\hspace*{\parindent}Die für die Witwenkassen-Angelegenheit ernannte Commission beehrt sich, das Resultat ihrer Berathung hiemit vorzulegen.

Es kam zuerst in Frage, welcher Zinsfuss als Grundlage für die zu treffenden Maassregeln anzunehmen sei.

Die Capitale der Witwenkasse sind zwar gegenwärtig dem grössten Theile nach noch zu höherm Zinsfuss als \(3\tfrac{1}{2}\) Procent belegt, jedoch so, dass bei richtiger Würdigung der Verhältnisse der mittlere Zinsfuss jedenfalls wie niedriger als 4 Procent stehend betrachtet werden muss. In Erwägung aber von folgenden Gründen:

1) dass ganz unverkennbar der Zinsfuss in allen Ländern Europas, einzelner Fluctuationen ungeachtet, die Tendenz zu weiterm Herabsinken zeigt\footnote{Einer der erleuchtetsten Staatsmänner, namentlich im Fache der Finanzverhältnisse, der Grossherz. Badensche Minister \textsc{Nebenius} sagt in seinem bekannten Werke Über die Herabsetzung der Zinsen der öffentlichen Schulden, S. 129: ‘Das allmählige weitere Sinken des Zinsfusses, welches bei längerer Fortdauer des Friedens nicht ausbleiben kann, wird zuletzt überall, hier etwas früher, dort etwas später, die Reduction auf drei Procent herbeiführen, oder sie wenigstens als eine nur von dem Entschlusse der Regierung abhängige Maassregel erscheinen lassen’; und an einer andern Stelle S. 21: ‘Einer Periode grösserer Regsamkeit in productiven Unternehmungen, die das Sinken des Zinsfusses eine Zeitlang aufhält, folgt um so gewisser ein rasches Sinken des Zinsfusses nach.’};

2) dass jede vom Zinsfusse wesentlich abhängige Anstalt, wenn sie nicht für eine durchaus unsichere gelten soll, nicht auf dem augenblicklich bestehenden, sondern auf einem etwas niedrigern Zinsfuss basirt werden muss;

3) dass auch die Regulirungen von 1794 uns in so fern mit einem guten Beispiele vorangegangen sind, als man damals wenigstens die Intention hatte, die Casse für den Zinsfuss von 3 Procent sicher zu stellen;

\noindent hält die Commission die Zugrundelegung eines höhern Zinsfusses als \(3\tfrac{1}{2}\) Proc. nicht für zulässig, und daher für unumgänglich nothwendig, dass man sich die Deckung des bei diesem Zinsfusse resultirenden Deficits von 17520 Thalern zum Ziele setze.

Wenn man vermeinte, dass die auch bei den bisherigen Einrichtungen zur Zeit noch Statt findenden jährlichen Überschüsse {\renewcommand{\thefootnote}{**)}\footnote{Die Jahresrechnung 1844—1845 ergibt keinen Überschuss; denn
\par\noindent für 1844 Juli 1 war das Geldvermögen ausgeworfen zu 116375 Thl. 14 Ggr. 6 Pf.
\par\noindent für 1845 Juli 1 aber, nur zu . . . . . . . . . . . . . . . . . 116369 ,, 17 ,, 4 ,,
\par\noindent Da jedoch im laufenden Jahre drei Witwen weniger sind, so wird in demselben wieder auf Überschuss zu rechnen sein, der auch wahrscheinlich noch eine Zeitlang fortdauern wird.}} schon an sich Schritte zur Deckung des Deficits seien, und dass dieses getilgt sein werde, sobald nur das jetzige Vermögen sich um 17520 Thaler vergrössert haben würde, so würde eine solche Meinung nur auf einer Begriffsverwirrung und auf einer gänzlichen Verkennung der Bedeutung des Deficits beruhen. Im Geiste des Calcüls, durch welchen die Grösse des Deficits eruirt ist, liegt in dem Resultate implicite schon die volle Berücksichtigung der jetzt noch Statt findenden jährlichen Ersparnisse mit, und diese wie Schritte zur Deckung betrachten, würde dasselbe sein, als wenn man sie zweimal in Rechnung brächte {\renewcommand{\thefootnote}{***)}\footnote{Da ähnliche Ansichten dem Vernehmen nach hie und da geäussert sind, in diesem Bericht aber zur Beseitigung derselben nicht der Ort sein würde, so ist der Anlage noch eine nähere Beleuchtung des Gegenstandes beigefügt, obwohl dieselbe für alle, welche die Denkschrift schon mit der nöthigen Aufmerksamkeit gelesen und erwogen haben, ganz überflüssig sein wird.}}. Die Genossenschaft möge sich also keine Illusion über die Wahrheit
